\subsection{Permutations without small cycles.}
%
%

Sharp bounds on $\PR(C_{[m]}(\sigma)=0)$ are a
key to establishing the Poisson model.
The model predicts that $\PR(C_{[m]}(\sigma)=0)$
should be about $\er^{-H_m}$,
and Corollary \ref{cycles-single-set} contains an upper bound
close to this.
 This cannot be expected to hold for large $m$, for example 
 $\PR(C_{[m]}(\sigma)=0)=1/n$ if $m \ge n/2$ since a permutation lacking cycles of length at most $m$ must be a single $n$-cycle.
  In fact, when $n/m$ is small, there is an asymptotic formula
 $\PR(C_{[m]}(\sigma)=0) \sim \omega(n/m)/m$ ($n\to\infty$, $m\to\infty$) where $\omega$ is Buchstab's function
 and $\omega(u)\to \er^{-\gamma}$ as $u\to\infty$ \cite[Theorem 5]{granville}.  This is analogous to the problem
 of counting integers $n\le x$ with no prime factor $\le x^{1/u}$
 (see \cite[Ch. III.6]{Tenbook}).

Our focus  is to prove that
$\PR(C_{[m]}(\sigma)=0)$ is very close to $\er^{-H_m}$ when 
$n/m$ is large.

 \begin{thm}\label{nosmallcycles2}
 Let $1\le m\le n$.  Then
 \[
 \PR(C_{[m]}(\sigma)=0) = \er^{-H_{m}} \(1 + O(\er^{-g(n/m)}) \),
 \]
 where $g(x)=0$ for $1\le x\le 20$ and for $x > 20$,
 \[
g(x) = x\log x - x\log\log\log x + O(x).
 \]
 \end{thm}
 
 Theorem \ref{nosmallcycles2} will be proved in section \ref{sec:nosmall}.
